% !TeX program = xelatex
% Template inicial para o artigo da disciplina PCS5917 -- IA Adversarial.
% Compilação inicial: xelatex artigo-jbcs. Após adicionar citações: xelatex, bibtex, xelatex e xelatex.
\documentclass{sbc2023}

\usepackage{graphicx}
\usepackage[misc,geometry]{ifsym}
\usepackage{fontspec}
\usepackage{fontawesome}
\usepackage{academicons}
\usepackage{color}
\usepackage{hyperref}
\usepackage{aas_macros}
\usepackage[bottom]{footmisc}
\usepackage{supertabular}
\usepackage{afterpage}
\usepackage{url}
\usepackage{pifont}
\usepackage{multicol}
\usepackage{multirow}


\setcitestyle{square}

\definecolor{orcidlogo}{rgb}{0.37,0.48,0.13}
\definecolor{unilogo}{rgb}{0.16,0.26,0.58}
\definecolor{maillogo}{rgb}{0.58,0.16,0.26}
\definecolor{darkblue}{rgb}{0.0,0.0,0.0}
\hypersetup{colorlinks,breaklinks,
  linkcolor=darkblue,urlcolor=darkblue,
  anchorcolor=darkblue,citecolor=darkblue}
% Descomente durante a redação se ocorrer erro de hyperlinks no PDF.
% \hypersetup{draft}

% Estes metadados serão definidos pela editoria na versão publicada.
\jid{JBCS}
\jtitle{Journal of the Brazilian Computer Society, 202X, XX:1,}
\doi{10.5753/jbcs.202X.XXXXXX}
\copyrightstatement{This work is licensed under a Creative Commons Attribution 4.0 International License}
\jyear{202X}

\title[Título curto]{Título completo do artigo}

% ORCID é obrigatório para cada autor no JBCS.
\author[Silva et al. 202X]{
\affil{\textbf{Caio Vinícius Ribeiro da Silva}~
  \href{https://orcid.org/0000-0000-0000-0000}{\textcolor{orcidlogo}{\aiOrcid}}~
  \textcolor{blue}{\faEnvelopeO}~~[
  \textbf{Universidade de São Paulo}~|\href{mailto:email@usp.br}{\textbf{\textit{email@usp.br}}}]}

% Copie o bloco \affil acima para cada coautor.
}

\renewcommand{\refname}{Referências}

\begin{document}

\begin{frontmatter}
\maketitle

\begin{mail}
Instituição, endereço completo, cidade, estado, CEP e país.
\end{mail}

\begin{dates}
% Informação fornecida pela editoria antes da publicação.
\small{\textbf{Recebido:} DD mês AAAA~~~$\bullet$~~~\textbf{Aceito:} DD mês AAAA~~~$\bullet$~~~\textbf{Publicado:} DD mês AAAA}
\end{dates}

\begin{abstract}
\textbf{Resumo.~}
Apresente o problema de pesquisa, o objetivo, o método, os principais resultados e a contribuição.
\end{abstract}

\begin{keywords}
Aprendizado de Máquina Adversarial, Modelos de Linguagem de Grande Porte, Segurança, Palavra-chave 4, Palavra-chave 5
\end{keywords}

\end{frontmatter}

\section{Introdução}
\label{sec:introduction}

Apresente o problema, sua relevância, a lacuna de pesquisa e as contribuições deste trabalho. Cite trabalhos relacionados com comandos como \verb|\citep{chave}|.

\section{Fundamentação e Trabalhos Relacionados}
\label{sec:background}

Apresente os conceitos necessários e posicione criticamente este trabalho em relação à literatura.

\section{Metodologia}
\label{sec:methodology}

Descreva o desenho da pesquisa, os conjuntos de dados, os modelos, os procedimentos, as métricas de avaliação e as condições de reprodutibilidade.

\subsection{Considerações Éticas}
\label{sec:ethics}

Descreva explicitamente como questões éticas, riscos, aprovações e o tratamento responsável dos dados foram considerados.

\section{Resultados e Discussão}
\label{sec:results}

Apresente os resultados e discuta suas implicações e limitações. Mantenha tabelas e figuras em \texttt{images/}.

% Exemplo de figura:
% \begin{figure}
%   \centering
%   \includegraphics[width=\columnwidth]{images/example.pdf}
%   \caption{Legenda que descreve a figura.}
%   \label{fig:example}
% \end{figure}

\section{Conclusão}
\label{sec:conclusion}

Resuma os achados, as limitações e os trabalhos futuros.

\section*{Declarações}

\begin{contributions}
Descreva a contribuição de cada autor, de preferência usando a Taxonomia CRediT (\href{https://credit.niso.org/}{https://credit.niso.org/}). Esta declaração é obrigatória.
\end{contributions}

\begin{interests}
Os autores declaram não possuir conflitos de interesse.
\end{interests}

\begin{acknowledgements}
Agradecimentos opcionais.
\end{acknowledgements}

\begin{funding}
Informe as fontes de financiamento, se houver.
\end{funding}

\begin{materials}
Informe onde os conjuntos de dados, o código, os prompts e outros materiais suplementares estão disponíveis; caso contrário, informe se poderão ser disponibilizados mediante solicitação. Esta declaração é obrigatória.
\end{materials}

\bibliographystyle{apalike-sol}
\bibliography{refs}

\end{document}
